\section{结论与未来工作}
\label{sec:conclusion}

在本研究中，我们构建了 \textbf{Grox\_chat}，一个基于黑板模式、由 LangGraph 有向无环图深度调度的异构多智能体推理论坛。本框架的核心哲学是：不要指望廉价模型拥有完美的工具调用能力与自律性，而是应该通过系统工程的手段，构建一个高度受控、强物理约束的“监狱式”实验场。

通过将原始推理与系统级宏观调度（如规划、终止判定、事实验证）解耦，引入动态排序的 DAPA 惩罚机制，以及创新的“自然语言提取意图”混合检索技术，Grox\_chat 成功地将一群不稳定、极易产生幻觉的廉价大语言模型，驯化为了一支能够持续进行高强度逻辑博弈的科研专家团队。

Grox\_chat 作为该架构的最底层（Chat-Only）基线模型，已经证明了通过严格的状态限制和异构分工，可以有效压榨出廉价模型的极致推理潜力，并大幅逼近昂贵单一巨型模型的分析质量。

\subsection{未来工作}
在当前的纯对话基准之上，我们计划沿着控制权逐步释放的路径，推出两个进阶版本：
\begin{enumerate}
    \item \textbf{Grox\_lite}：引入受限的代码生成能力。在这一版本中，智能体将被允许生成用于数据分析与数学验证的 Python 代码片段，系统底层会接管并隔离执行这些代码，将其输出作为硬事实返回给黑板。
    \item \textbf{Grox\_pro}：实现全自动的“自主科学实验室”。最高级别的系统将打通与 OpenClaw 级别的外部沙盒环境交互，允许模型调用真实的物理工具（如远程 SSH、容器内 CLI 执行），从而实现从“提出假设”到“运行验证”的完整科研闭环。
\end{enumerate}